% Tabasco REPA variant grid -- \input{} from report-draft.tex Ch5.
\begin{table}[tbp]
\centering
\small
\begin{tabular}{l l c l l}
\toprule
\textbf{Model} & \textbf{Encoder} & \textbf{Enc.\ dim} & \textbf{Combination} & \textbf{Plumbing} \\
\midrule
Baseline & --- & --- & --- & --- \\
\midrule
CheMeleon additive (same)    & CheMeleon (2D) & 2048 & additive & separate \\
CheMeleon additive (fused)   & CheMeleon (2D) & 2048 & additive & fused \\
CheMeleon tradeoff (same)    & CheMeleon (2D) & 2048 & tradeoff & separate \\
CheMeleon tradeoff (fused)   & CheMeleon (2D) & 2048 & tradeoff & fused \\
MACE additive                & MACE (3D)      & 192  & additive & fused \\
MACE tradeoff                & MACE (3D)      & 192  & tradeoff & fused \\
\bottomrule
\end{tabular}
\caption{The REPA variant grid for Tabasco. \emph{Combination} sets how the REPA term enters the loss (\emph{additive}: a plain regulariser added to the generative objective; \emph{tradeoff}: a convex combination that down-weights the generative term). \emph{Plumbing} sets how the trunk's coordinate and atom-type hidden states reach the projector (\emph{fused}: concatenated into a single aligned representation; \emph{separate}: two projections). The grid is not a clean cube: both axes were swept for CheMeleon, while MACE was run only in the fused configuration. Our reading centres on additive combination with fused plumbing.}
\label{tab:tabasco-variants}
\end{table}
